\bookchapter{第六十七章：打磨与发布——测试、调试与上架 App Store}

\chapterintro{App 做出来了，然后呢？这一章回答“从能跑”到“能交付”的三件事：怎么高效地调试找 bug、怎么写测试防回归、怎么把 App 送上 App Store 让全世界下载。上架不只是仪式感——它是把一个作品变成产品的分水岭。}

\booksection{一、调试三板斧：print、断点、预览}

\begin{center}
\begin{tabularx}{\linewidth}{|X|X|X|}
\hline
\textbf{工具} & \textbf{干什么} & \textbf{怎么用} \tabularnewline
\hline
print + 控制台 & 看运行日志 & 代码里 print("值=\textbackslash{}(x)")，Xcode 底部控制台看输出 \tabularnewline
断点调试 & 让程序停在某行，看变量值 & 行号旁点一下出蓝标，运行到那里暂停，逐步走 \tabularnewline
SwiftUI 预览 + 调试器 & 实时看界面与视图层级 & \#Preview 画布 + 运行时的 View Hierarchy 检查器 \tabularnewline
\hline
\end{tabularx}
\end{center}

\begin{center}\small\color{gray}调试三板斧\end{center}

\begin{tcbinfo}{小贴士：崩溃了怎么读？}
Swift 崩溃时控制台会打印 “Fatal error” 加原因，Xcode 会停在出错那一行并高亮。真机上的历史崩溃可以在 Xcode Window Organizer Crashes 里看（需付费账号）。关键词优先搜\textbf{“错误原因第一句 + swiftui”}——精确的问题，才有精确的答案。

\end{tcbinfo}

\booksection{二、测试：给代码上保险}

手测靠不住——改了这个功能，另一个悄悄坏了。测试就是“后悔药”兼“防弹衣”。iOS 用\textbf{XCTest}框架，分两种：

\begin{center}
\begin{tabularx}{\linewidth}{|X|X|X|X|}
\hline
\textbf{层级} & \textbf{测什么} & \textbf{Target} & \textbf{速度} \tabularnewline
\hline
单元测试 & 纯逻辑（计算、校验） & MyFirstAppTests & 秒级 \tabularnewline
UI 测试 & 真实界面行为 & MyFirstAppUITests & 分钟级 \tabularnewline
手动测试 & 真人感受、体验 & —— & 永远需要 \tabularnewline
\hline
\end{tabularx}
\end{center}

\begin{center}\small\color{gray}iOS 测试层级（模板项目已预置两个 Test Target）\end{center}

\begin{lstlisting}[language=Swift]
import XCTest

// 把"金额字符串转数字"这种逻辑抽成纯函数
func parseAmount(_ input: String) -> Double? {
    guard let v = Double(input.trimmingCharacters(in: .whitespaces)),
          v > 0 else { return nil }
    return v
}

final class AmountTests: XCTestCase {
    func testEmptyInputReturnsNil() {
        XCTAssertNil(parseAmount(""))
    }
    func testNegativeReturnsNil() {
        XCTAssertNil(parseAmount("-5"))
    }
    func testNormalValueParses() {
        XCTAssertEqual(parseAmount("12.5"), 12.5)
    }
}
\end{lstlisting}

跑测试：CmdU。全绿之后再改代码，心里有底。记账本里的金额校验、汇总计算，都该有单元测试罩着。

\booksection{三、签名与打包：给 App 盖上“印章”}

每个 iOS App 都必须\textbf{签名}才能安装——证明“这个包出自你，没被篡改”。层级关系：

\begin{itemize}
  \item \textbf{学习阶段}：免费 Apple ID + Xcode 自动签名（Personal Team），真机可跑，7 天续期一次；
  \item \textbf{分发阶段}：付费开发者账号（99 美元/年）签发证书与描述文件，才能上架和长期安装。
\end{itemize}

打包：Xcode 菜单 Product Archive（归档），构建完成后在 Organizer 里Distribute App选 App Store Connect 上传。\textbf{归档必须用 Generic iOS Device 作为目标}（不是模拟器），这是新手常忘的一步。

\textbf{避坑提醒：证书丢了 = 更新断档}\par

签名证书和私钥一旦丢失，你无法再给旧 App 发更新——用户只能装“全新”应用。妥善备份证书（钥匙串导出 .p12）并记好密码。好在苹果提供了“自动签名”和云证书管理（Xcode 里勾选 Automatically manage signing），能省掉大部分手工操作的坑。

\booksection{四、上架 App Store 全流程}

注册开发者账号

准备商店素材

上传构建包

填隐私与审核信息

苹果审核（1-7 天）

上架

\begin{center}\small\color{gray}图 5：App Store 上架六步走\end{center}

\begin{itemize}
  \item \textbf{注册}：developer.apple.com，个人账号 99 美元/年；
  \item \textbf{素材}：图标（1024×1024，别带圆角——苹果自己切）、截图（iPhone 各尺寸）、描述、隐私政策链接；
  \item \textbf{隐私问卷}：如实申报收集了哪些数据——不联网的记账本可以"不收集任何数据"；
  \item \textbf{审核}：首次提交常有 1-7 天等待，被拒就按"审核备注"逐条改——苹果的拒绝理由通常写得非常具体，照做即可；
  \item \textbf{TestFlight}：上架前先把测试版发给朋友内测（最多 10000 人），收集一轮反馈再正式提交，专业范儿十足。
\end{itemize}

\booksection{五、商店素材怎么做：第一印象就是下载量}

商店页是用户“逛店”时看到的橱窗。三个素材决定第一印象：

\begin{itemize}
  \item \textbf{图标}：1024×1024，别用截图凑合。没设计能力就用 AI 工具生成一个简洁图形，\textbf{在桌面上缩小后还能一眼认出}才算合格；
  \item \textbf{截图}：至少 2 套（iPhone 主流尺寸），展示最核心的功能场景；每张配一句人话文案（"一秒记一笔"比"高效的记账解决方案"好十倍）；
  \item \textbf{描述}：第一句说清"这是什么 + 给谁用"，别堆砌"最专业、最强"——真诚比夸张有效。
\end{itemize}

\begin{tcbpractice}{来动手：把记账本"交付"一次}
\booksection{六、崩溃上报与性能优化小清单}

App 上架前，这两件事值得各花半小时：

\booksubsection{崩溃监控}

你测试时没崩，不代表用户手机不崩。苹果在 Xcode Organizer 里提供免费的崩溃报告（付费账号可看真机崩溃日志）；更省心的方案是接入第三方 SDK（如 Firebase Crashlytics）。\textbf{没有崩溃监控的上架，等于闭着眼睛开车}。

\booksubsection{性能自查四条}

\begin{enumerate}
  \item \textbf{列表用 LazyVStack / List}：长列表别用普通 VStack 硬排——懒加载是性能第一课；
  \item \textbf{图片别硬塞原图}：用 AsyncImage（远程图）和系统缓存，5MB 原图足以卡爆老机型；
  \item \textbf{别在主线程干重活}：网络、大文件都走 async/await；
  \item \textbf{用 Instruments 量一量}：Xcode 自带的性能分析器（Time Profiler、内存泄漏检查）——先测量，再优化，别凭感觉。
\end{enumerate}

(1) 给 parseAmount 等逻辑补 3 个单元测试（CmdU 跑绿）；(2) 用断点亲手抓一个自己制造的崩溃并修好；(3) Product Archive 打一个包，走 TestFlight 邀请朋友内测；(4) 有余力就注册开发者账号，把记账本真正提交上架。前三步今晚就能做完。

\end{tcbpractice}

\begin{tcbsummary}{本章小结}
\begin{itemize}
  \item 调试三板斧：print 日志、断点、SwiftUI 预览与视图检查器
  \item XCTest：单元测试（CmdU）+ UI 测试 + 手动测试，模板已预置
  \item 签名：免费账号自动签名（7 天续期）；付费账号才能分发
  \item 打包：Generic iOS Device + Archive + Distribute；证书务必备份
  \item 上架六步：注册 素材 上传 隐私问卷 审核 上架；TestFlight 先内测
\end{itemize}

\end{tcbsummary}
